% ============================================================
% litreview.tex
% ============================================================

This chapter surveys three bodies of literature that together
motivate and situate our empirical strategy: the statistical
foundations of survival analysis and competing-risks methods;
the empirical record on securities class action outcomes,
duration, and procedural reform; and the structure and known
limitations of the FJC IDB. We conclude by articulating the specific gap that
this thesis fills.

\section{Survival Analysis and Competing Risks:
Conceptual Foundations}

Survival analysis models time-to-event data in the presence of
censoring \citep{kalbfleisch2002}. The Kaplan-Meier estimator
provides nonparametric survival curves, while the Cox
proportional hazards model estimates semi-parametric covariate
associations without specifying the baseline hazard
\citep{kaplan1958, cox1972}. In settings where cases face
multiple mutually exclusive exits, however, standard
single-event survival methods are insufficient.

That problem is central here. Securities class actions
typically terminate through either settlement or dismissal,
and the occurrence of one precludes the other. Treating one
outcome as the event of interest and the other as ordinary
censoring biases cumulative incidence estimates
\citep{putter2007}. The competing-risks literature instead
distinguishes between cause-specific hazard models, which
estimate event-specific instantaneous rates, and Fine-Gray
subdistribution models, which estimate how covariates shift
the cumulative incidence function (CIF) itself
\citep{fine1999}. The Aalen-Johansen estimator provides the
nonparametric analogue for CIF estimation, and Gray's test
compares CIFs across groups \citep{aalen1978, gray1988}.

While survival and competing-risks methods are standard in
clinical trials and labor economics, they remain uncommon in
empirical studies of civil litigation. Exceptions include work
on bankruptcy case durations \citep{dinterman2021,
dewaelheyns2009}. Applying these methods to securities
litigation raises two additional methodological challenges:
separating the effect of a legislative reform from concurrent
changes in case composition, and accounting for unobserved
heterogeneity across judicial circuits.

\subsection{Composition Adjustment and Unobserved Heterogeneity}

Two methodological extensions augment the core competing-risks
framework in this thesis: propensity-score re-weighting for
composition adjustment, and shared frailty models for
unobserved clustering.

Inverse probability of treatment weighting uses a propensity
score to re-weight observations so that treated and control
groups are balanced on observed covariates
\citep{xieliu2005}. \citet{austinfine2025} extend IPTW to
competing-risks settings, showing that weighted Aalen-Johansen
and cause-specific Cox estimators can recover adjusted
cumulative incidence and hazard estimates under strong
identifying assumptions. In our setting, those assumptions are
not cleanly satisfied, so IPTW is used as a composition
adjustment rather than as a causal design; the precise scope
of that interpretation is developed in
Chapter~\ref{ch:methodology}.

Shared frailty models address a different source of bias:
unobserved heterogeneity within clusters. When cases filed
in the same judicial circuit share latent characteristics---local
judicial norms, bench expertise, or litigation culture---that
are not captured by observed covariates, standard Cox
estimates that pool across circuits may be biased
\citep{ruetenbudde2019}. A shared frailty model introduces
a circuit-level random effect (the ``frailty'') that captures
this unobserved clustering. \citet{ruetenbudde2019} develop
competing-risks frailty models for multicenter clinical data,
where center-level frailties capture institutional
heterogeneity in baseline risk---directly analogous to
modeling unobserved jurisdictional variation across federal
circuits. \citet{frevent2024} extend this framework to
spatially correlated frailties, treating geographic units
as clusters with shared latent risk factors.

The small number of circuit-level clusters in our setting
constrains the precision of frailty variance estimates; the
specific implications and our mitigation strategy are
addressed in Chapter~\ref{ch:methodology}.

\section{Empirical Landscape of Securities Class Actions}

\subsection{Outcome Frequencies and Resolution Pathways}

Practitioner reports and academic datasets consistently show
that securities class actions almost never proceed to trial.
Cornerstone Research reports that approximately 46\% of resolved
cases settle and roughly 43\% are dismissed, with fewer than
1\% reaching trial \citep{cornerstone2022}. Dismissals tend to
occur earlier than settlements: about half of all dismissals
occur within three years of filing, whereas median times to
settlement range from 3.0 to 3.7 years \citep{cornerstone2024}.
These descriptive cohort plots provide valuable baseline
information but do not estimate formal hazard functions, account
for censoring, or test for covariate effects while properly
modeling the competing nature of outcomes.

\subsection{The PSLRA and Procedural Reform}

The PSLRA represents
the most significant structural break in the history of
securities class action litigation. The Act introduced
heightened pleading standards requiring plaintiffs to plead
facts giving rise to a ``strong inference'' of scienter, imposed
an automatic stay of discovery pending resolution of motions to
dismiss, and established lead-plaintiff provisions favoring
institutional investors \citep{johnson2001}. Empirical studies
have found that the PSLRA was associated with higher early
dismissal rates for cases lacking strong merit signals---such as accounting
restatements, parallel SEC enforcement actions, or insider
trading allegations---while allowing meritorious cases to
proceed \citep{johnson2001, choi2007}.

However, existing studies have not formally tested whether PSLRA 
effects vary over litigation duration, whether they differ between
cause-specific and subdistribution interpretations, or whether
the reform's impact was heterogeneous across judicial circuits.
Consequently, it remains an open question in the literature 
whether the PSLRA operated uniformly as a static filter, or 
whether its procedural mechanisms generated time-dependent and 
geographically fragmented impacts on case resolution.

\subsection{Determinants of Settlement and Duration}

Academic research has identified several factors associated
with settlement amounts and case duration. The magnitude of
alleged investor losses is the strongest predictor of
settlement size \citep{cox2008}. Larger cases are more likely
to survive dismissal and settle for substantial amounts, but
they also take longer to resolve \citep{cornerstone2024}.
Cases with institutional lead plaintiffs tend to settle for
higher amounts but exhibit longer durations
\citep{cheng2010, cox2008}.

Other determinants include merit signals (accounting
restatements, SEC investigations), statutory basis
(Section~11 versus Section~10(b) claims), and forum variation
across federal circuits \citep{johnson2001, perino2003}.
\citet{wang2022} find that law firm market share is associated
with case outcomes and settlement patterns.
Furthermore, multidistrict litigation (MDL) consolidation represents a
distinct procedural pathway that concentrates large numbers
of related cases before a single judge. Yet, the current
literature lacks a dynamic analysis of how this consolidation
alters the underlying hazard processes of case resolution
over time. Because MDL consolidation suppresses both
settlement and dismissal rates simultaneously, its net
effect on cumulative outcome probabilities cannot be
determined from static models that examine each outcome
in isolation.

\subsection{Prior Use of Survival Methods}

While survival methods are rare in this domain, \citet{brochet2014}
analyze settlement speed using Cox proportional hazards models
conditional on cases that eventually settle. That design
mischaracterizes the data-generating process. Because
dismissals tend to occur earlier than settlements, treating
dismissed cases as right-censored implies that they remain at
risk of settling at some unobserved future time, which
upwardly biases settlement probabilities and distorts hazard
ratios. The limitation of this conditional-on-settlement
approach motivates the competing-risks framework developed in
this thesis.

\section{Federal Court Administrative Data}

Our analysis relies on the IDB, which aggregates civil case data reported by
the Administrative Office of the U.S.\ Courts from all 94
federal district courts \citep{fjc_idb_portal, fjc_idb_guide}.
The IDB contains longitudinal records from approximately 1970
to the present, capturing filing dates, termination dates,
nature-of-suit codes, disposition codes, and court identifiers.
Securities class actions are primarily identified using
nature-of-suit (NOS) Code~850 (``Securities, Commodities,
and Exchange'').

A major challenge is outcome coding. Voluntary dismissals may
include settlement-like resolutions, and judgment-bearing
dispositions can conflate plaintiff and defendant victories
unless the IDB's \texttt{JUDGMENT} field is consulted.
Docket-based work and prior audits therefore suggest that raw
IDB disposition fields can misstate settlement counts and
misclassify outcomes \citep{scales2024, clermont2002,
hadfield2004}. These measurement problems make the IDB
analytically valuable but not turnkey; the dedicated coding
framework is developed in Chapter~\ref{ch:data}.

Despite these limitations, the IDB remains the most
comprehensive source for population-level data on federal civil
litigation and has been widely used in empirical legal research
\citep{clermont2002, scales2024}.

\section{Research Gap and Contribution}

The literature establishes that (1) survival and
competing-risks methods are well-suited to modeling
time-to-event data with multiple outcome types; (2) securities
class actions exhibit predictable patterns in outcome
frequencies and timing, with the PSLRA as a major structural
break; and (3) the IDB provides rich administrative data but
requires careful handling of disposition codes and class-action
flags.

Despite these foundations, existing empirical work has not
estimated the specific objects that a competing-risks framework
provides. Practitioner reports provide descriptive cohort plots
without formal hazard modeling \citep{cornerstone2024}.
Academic studies use Cox models conditional on settlement
without accounting for competing risks \citep{brochet2014},
or analyze outcomes using static regression methods that ignore
duration \citep{cox2008, johnson2001}. None estimate joint
cause-specific hazard functions, subdistribution hazards, and
cumulative incidence curves for settlement and dismissal while
properly accounting for censoring and the mutually exclusive
nature of outcomes. Furthermore, no prior study has attempted
to disentangle the PSLRA's legislative effect from concurrent
changes in case composition using propensity-score methods, nor
has any modeled unobserved circuit-level heterogeneity through
shared frailty specifications in this litigation context.

This thesis fills the gap by applying a full
competing-risks framework---cause-specific and
subdistribution hazard models, propensity-score composition
adjustment, and shared frailty sensitivity
analysis---to population-level federal securities litigation
data. Where prior work estimates single-outcome hazards or
static regression coefficients, this thesis jointly
estimates cumulative incidence, time-varying and
circuit-heterogeneous PSLRA associations, and the distinct
contributions of case composition versus the legislative
reform itself. The following chapter formalizes the
statistical framework used to pursue these objectives.
